\chapter{Εισαγωγή}
\label{chap:introduction}

\section{Διατύπωση προβλήματος}

\subsection{Ο κτιριακός τομέας στο ενεργειακό και κλιματικό πρόβλημα}

Η ενεργειακή μετάβαση δεν είναι πλέον ένα αφηρημένο ζήτημα ενεργειακής
πολιτικής, αλλά ένα πρακτικό πρόβλημα που επηρεάζει το κόστος διαβίωσης, την
ποιότητα των κατοικιών και τον τρόπο με τον οποίο σχεδιάζονται οι δημόσιες
παρεμβάσεις. Η ανάγκη περιορισμού των εκπομπών αερίων του θερμοκηπίου, η
αστάθεια των αγορών ενέργειας, η εξάρτηση από εισαγόμενα καύσιμα και η αύξηση
του κόστους θέρμανσης και ψύξης έχουν στρέψει μεγάλο μέρος της συζήτησης στην
εξοικονόμηση ενέργειας. Μέσα σε αυτή τη συζήτηση, ο κτιριακός τομέας έχει
ιδιαίτερη σημασία, επειδή συνδέεται άμεσα με την καθημερινή κατανάλωση
ενέργειας των νοικοκυριών, των επιχειρήσεων και του δημόσιου τομέα.

Σε παγκόσμιο επίπεδο, ο Διεθνής Οργανισμός Ενέργειας
(\eng{International Energy Agency} -- IEA) εκτιμά ότι η
λειτουργική χρήση ενέργειας στα κτίρια αντιστοιχεί περίπου στο 30\% της
παγκόσμιας τελικής κατανάλωσης ενέργειας, ενώ το ποσοστό αυξάνεται όταν
συνυπολογιστεί η ενέργεια που σχετίζεται με βασικά κατασκευαστικά υλικά
\cite{iea_buildings_2026}. Αντίστοιχα, η έκτη έκθεση αξιολόγησης του
\eng{Intergovernmental Panel on Climate Change} αντιμετωπίζει τα κτίρια ως
διακριτό πεδίο μετριασμού \cite{ipcc_ar6_wg3_buildings}. Σε ευρωπαϊκό επίπεδο,
η Ευρωπαϊκή Επιτροπή αναφέρει ότι περίπου το 40\% της ενέργειας που
καταναλώνεται στην Ευρωπαϊκή Ένωση σχετίζεται με τα κτίρια, ενώ μεγάλο μέρος
του κτιριακού αποθέματος παραμένει ενεργειακά ανεπαρκές
\cite{ec_epbd_page_2026,ec_epbd_news_2025}. Τα βασικά μεγέθη που τεκμηριώνουν
την κλίμακα του προβλήματος συνοψίζονται στο Σχήμα~\ref{fig:buildings-energy-context}.

\begin{figure}[!htbp]
    \centering
    \begin{adjustbox}{max width=\textwidth,center}
        \begin{tikzpicture}[
    card/.style={
        rectangle,
        rounded corners=3pt,
        draw=blue!55,
        fill=blue!6,
        align=center,
        text width=5.35cm,
        minimum height=1.65cm,
        font=\sffamily\small
    }
]
\node[card] (global) at (-3.0,1.05) {
    {\Large\bfseries\textcolor{blue!65!black}{30\%}}\\[3pt]
    Παγκόσμια τελική κατανάλωση\\
    ενέργειας στα κτίρια
};

\node[card] (euenergy) at (3.0,1.05) {
    {\Large\bfseries\textcolor{blue!65!black}{40\%}}\\[3pt]
    Κατανάλωση ενέργειας στην ΕΕ\\
    που σχετίζεται με κτίρια
};

\node[card] (oldstock) at (-3.0,-1.05) {
    {\Large\bfseries\textcolor{blue!65!black}{85\%}}\\[3pt]
    Κτίρια της ΕΕ κατασκευασμένα\\
    πριν από το 2000
};

\node[card] (inefficient) at (3.0,-1.05) {
    {\Large\bfseries\textcolor{blue!65!black}{75\%}}\\[3pt]
    Κτιριακό απόθεμα χαμηλής\\
    ενεργειακής απόδοσης στην ΕΕ
};
\end{tikzpicture}

    \end{adjustbox}
    \caption[Ενδεικτικοί δείκτες για τον κτιριακό τομέα και την ανάγκη ανακαινίσεων]{Ενδεικτικοί δείκτες που αναδεικνύουν τη σημασία του κτιριακού τομέα στην ενεργειακή κατανάλωση και στην ανάγκη ανακαινίσεων. Οι δείκτες έχουν διαφορετικό παρονομαστή και παρουσιάζονται ως συνοπτική εικόνα του προβλήματος. Ιδία επεξεργασία βάσει στοιχείων IEA και European Commission \cite{iea_buildings_2026,ec_epbd_page_2026,ec_epbd_news_2025}.}
    \label{fig:buildings-energy-context}
\end{figure}

Η κατανάλωση ενέργειας στα κτίρια προκύπτει κυρίως από θέρμανση, ψύξη, ζεστό
νερό χρήσης, φωτισμό και ηλεκτρικές συσκευές. Ιδίως στις κατοικίες, η ζήτηση
ενέργειας επηρεάζεται από την ηλικία του κτιρίου, την ποιότητα του κτιριακού
κελύφους, την απόδοση των τεχνικών συστημάτων και τον τρόπο χρήσης. Ως
κτιριακό κέλυφος νοούνται τα δομικά στοιχεία που οριοθετούν τον εσωτερικό
χώρο, όπως εξωτερικοί τοίχοι, οροφή, δάπεδο και κουφώματα, μαζί με το επίπεδο
θερμομόνωσής τους. Η Eurostat αναφέρει ότι το 2024 τα ιδιωτικά νοικοκυριά
κατανάλωσαν το
27\% της τελικής ενέργειας στην Ευρωπαϊκή Ένωση, ενώ η θέρμανση χώρων
αντιστοιχούσε περίπου στα δύο τρίτα αυτής της κατανάλωσης
\cite{eurostat_energy_2026}. Τα στοιχεία αυτά δεν αποδεικνύουν από μόνα τους
την επίδραση συγκεκριμένων ανακαινίσεων. Δείχνουν όμως ότι οι χρήσεις που
επηρεάζονται από παρεμβάσεις στο κτιριακό κέλυφος και στα τεχνικά συστήματα
αποτελούν κρίσιμο μέρος της ενεργειακής ζήτησης.

\subsection{Παλαιότητα και ενεργειακή ανεπάρκεια του κτιριακού αποθέματος}

Η ανάγκη για αναβάθμιση των κτιρίων συνδέεται άμεσα με τη σύνθεση του
υφιστάμενου κτιριακού αποθέματος. Στην Ευρωπαϊκή Ένωση, μεγάλο μέρος των
κτιρίων έχει κατασκευαστεί πριν από την εφαρμογή αυστηρών ενεργειακών απαιτήσεων.
Η Ευρωπαϊκή Επιτροπή έχει επισημάνει ότι το 85\% των κτιρίων της ΕΕ
κατασκευάστηκε πριν από το 2000 και ότι περίπου το 75\% του κτιριακού
αποθέματος έχει χαμηλή ενεργειακή απόδοση \cite{ec_epbd_news_2025}. Το στοιχείο
αυτό έχει ιδιαίτερη βαρύτητα, επειδή τα κτίρια έχουν μεγάλο κύκλο ζωής:
σημαντικό μέρος των σημερινών κτιρίων αναμένεται να χρησιμοποιείται και τις
επόμενες δεκαετίες. Συνεπώς, οι στόχοι κλιματικής ουδετερότητας δεν μπορούν
να επιτευχθούν μόνο μέσω νέων κτιρίων υψηλής απόδοσης, αλλά απαιτούν
συστηματική αναβάθμιση των υφιστάμενων.

Η ενεργειακή αναβάθμιση δεν αντιστοιχεί συνήθως σε μία μεμονωμένη τεχνική
επιλογή. Στην πράξη, συνδυάζει παρεμβάσεις στο κτιριακό κέλυφος, στα τεχνικά
συστήματα και, σε ορισμένες περιπτώσεις, στην τοπική παραγωγή ενέργειας. Επομένως, κάθε
ανακαίνιση είναι ταυτόχρονα τεχνικό έργο και επενδυτική απόφαση: απαιτεί
εκτίμηση κόστους, επιλογή παρεμβάσεων, πιθανή αναζήτηση χρηματοδότησης και
τεκμηρίωση του αναμενόμενου ενεργειακού αποτελέσματος.

\subsection{Ευρωπαϊκό θεσμικό πλαίσιο για την ενεργειακή απόδοση κτιρίων}

Η Ευρωπαϊκή Ένωση έχει αναπτύξει πολυεπίπεδο θεσμικό πλαίσιο για την
ενεργειακή απόδοση των κτιρίων. Η Ευρωπαϊκή Πράσινη Συμφωνία έθεσε την
κλιματική ουδετερότητα στο κέντρο της στρατηγικής της Ένωσης
\cite{ec_green_deal_2019}, ενώ η αναθεωρημένη Οδηγία (ΕΕ) 2024/1275 για την
ενεργειακή απόδοση των κτιρίων προβλέπει σταδιακή μετάβαση σε κτιριακό απόθεμα
υψηλής ενεργειακής απόδοσης και μηδενικών εκπομπών έως το 2050
\cite{eurlex_epbd_2024}. Η Οδηγία (ΕΕ) 2023/1791 για την ενεργειακή απόδοση
συμπληρώνει το πλαίσιο με γενικότερους στόχους εξοικονόμησης
\cite{eurlex_eed_2023}, ενώ τα εθνικά σχέδια ανακαίνισης κτιρίων αποτελούν το
εργαλείο με το οποίο τα κράτη μέλη καλούνται να εξειδικεύσουν την πορεία
αναβάθμισης του αποθέματός τους \cite{ec_nbrp_2026}. Η χρονική εξέλιξη των
βασικών ευρωπαϊκών πολιτικών αποτυπώνεται συνοπτικά στο
Σχήμα~\ref{fig:policy-timeline}.

\begin{figure}[!htbp]
    \centering
    \begin{adjustbox}{max width=\textwidth,center}
        \begin{tikzpicture}[
    year/.style={
        rectangle,
        rounded corners=2pt,
        draw=blue!65,
        fill=blue!65,
        text=white,
        align=center,
        minimum width=1.35cm,
        minimum height=0.62cm,
        font=\sffamily\bfseries\small
    },
    event/.style={
        rectangle,
        rounded corners=3pt,
        draw=blue!50,
        fill=blue!6,
        align=left,
        text width=9.6cm,
        minimum height=0.78cm,
        inner xsep=7pt,
        font=\sffamily\small
    },
    arrow/.style={-{Latex[length=2mm]}, thick, draw=gray!65}
]
\node[year] (y2010) at (0,0) {2010};
\node[event, right=0.45cm of y2010] (e2010) {EPBD 2010/31/ΕΕ: βασικό πλαίσιο για την ενεργειακή απόδοση των κτιρίων};

\node[year, below=0.55cm of y2010] (y2018) {2018};
\node[event, right=0.45cm of y2018] (e2018) {Τροποποίηση EPBD και ενίσχυση των μακροχρόνιων στρατηγικών ανακαίνισης};

\node[year, below=0.55cm of y2018] (y2020) {2020};
\node[event, right=0.45cm of y2020] (e2020) {Renovation Wave: επιτάχυνση και τουλάχιστον διπλασιασμός του ρυθμού ανακαινίσεων};

\node[year, below=0.55cm of y2020] (y2024) {2024};
\node[event, right=0.45cm of y2024] (e2024) {EPBD recast, Οδηγία (ΕΕ) 2024/1275: πορεία προς κτίρια μηδενικών εκπομπών};

\node[year, below=0.55cm of y2024] (y2050) {2050};
\node[event, right=0.45cm of y2050] (e2050) {Στόχος για υψηλής ενεργειακής απόδοσης και απανθρακοποιημένο κτιριακό απόθεμα};

\draw[arrow] (y2010) -- (y2018);
\draw[arrow] (y2018) -- (y2020);
\draw[arrow] (y2020) -- (y2024);
\draw[arrow] (y2024) -- (y2050);
\end{tikzpicture}

    \end{adjustbox}
    \caption[Εξέλιξη του ευρωπαϊκού πλαισίου για την ενεργειακή απόδοση κτιρίων]{Ενδεικτική εξέλιξη του ευρωπαϊκού πλαισίου για την ενεργειακή απόδοση κτιρίων. Ιδία επεξεργασία βάσει European Commission και EUR-Lex \cite{ec_renovation_wave_2026,eurlex_epbd_2024,ec_nbrp_2026}.}
    \label{fig:policy-timeline}
\end{figure}

Η στρατηγική \eng{Renovation Wave} εξειδίκευσε αυτή την κατεύθυνση, θέτοντας
ως στόχο την ανακαίνιση 35 εκατομμυρίων κτιρίων έως το 2030 και τον τουλάχιστον
διπλασιασμό του ετήσιου ρυθμού ενεργειακών ανακαινίσεων
\cite{ec_renovation_wave_2026,ec_renovation_wave_com_2020}. Το ευρωπαϊκό
πλαίσιο, επομένως, αντιμετωπίζει τις ανακαινίσεις όχι ως μεμονωμένες τεχνικές
παρεμβάσεις, αλλά ως διαρκή πολιτική που χρειάζεται κανόνες, χρηματοδοτικούς
μηχανισμούς και εργαλεία πληροφόρησης.

\subsection{Η ελληνική πραγματικότητα και τα προγράμματα ενεργειακής αναβάθμισης}

Στην Ελλάδα, η ενεργειακή αναβάθμιση κατοικιών συνδέεται τόσο με την κατάσταση
του κτιριακού αποθέματος όσο και με την οικονομική δυνατότητα των νοικοκυριών.
Η Απογραφή Κτιρίων της ΕΛΣΤΑΤ παρέχει στοιχεία για την περίοδο κατασκευής, τη
χρήση και άλλα χαρακτηριστικά των κτιρίων \cite{elstat_buildings_census_2021},
ενώ η Eurostat καταγράφει ότι το 2023 το ποσοστό πληθυσμού που ζούσε σε
κατοικία με πρόσφατη βελτίωση ενεργειακής απόδοσης ήταν στην Ελλάδα χαμηλότερο
από τον μέσο όρο της ΕΕ και ότι σημαντικό μέρος του πληθυσμού δήλωνε αδυναμία
ικανοποιητικής θέρμανσης τον χειμώνα
\cite{eurostat_living_conditions_energy_efficiency_2026}. Τα στοιχεία αυτά δεν
περιγράφουν ένα συγκεκριμένο πρόγραμμα, αλλά εξηγούν γιατί η αναβάθμιση
κατοικιών αποτελεί πρακτικό και κοινωνικό ζήτημα.

Το ελληνικό πλαίσιο διαθέτει επίσης μηχανισμούς τεκμηρίωσης της ενεργειακής
κατάστασης. Το Πιστοποιητικό Ενεργειακής Απόδοσης (ΠΕΑ) είναι υποχρεωτικό σε
περιπτώσεις πώλησης και ενοικίασης κτιρίων ή κτιριακών μονάδων, ενώ η
ενεργειακή κατηγορία πρέπει να δηλώνεται στις σχετικές εμπορικές καταχωρήσεις
\cite{ypen_kenak_2026}. Παράλληλα, η μακροπρόθεσμη στρατηγική ανακαίνισης
κτιρίων του Υπουργείου Περιβάλλοντος και Ενέργειας εντάσσει τις παρεμβάσεις
στον ευρύτερο στόχο σταδιακού μετασχηματισμού του κτιριακού αποθέματος
\cite{ypen_ltrs_2026}.

Σε επίπεδο χρηματοδότησης, η πιο γνωστή οικογένεια δημόσιων προγραμμάτων είναι
τα προγράμματα «Εξοικονομώ». Το «Εξοικονομώ 2025» συνεχίζει προηγούμενα σχήματα
και στοχεύει σε εξοικονόμηση πρωτογενούς ενέργειας άνω του 30\% για κάθε
κατοικία και αναβάθμιση κατά τουλάχιστον τρεις ενεργειακές κατηγορίες
\cite{exoikonomo2025_programme,exoikonomo2025_home,greece20_exoikonomo2025}.
Εκτός από τα δημόσια προγράμματα, υπάρχουν τραπεζικά προϊόντα που συνδέονται
με ενεργειακές παρεμβάσεις, πράσινα δάνεια ή συμπληρωματική χρηματοδότηση
\cite{eurobank_exoikonomo2025_2026,nbg_green_loan_2026}. Οι βασικές κατηγορίες
χρηματοδοτικών εργαλείων που ενδιαφέρουν την παρούσα εργασία συνοψίζονται στον
Πίνακα~\ref{tab:financing-categories}.

\begin{table}[!htbp]
    \centering
    \caption[Κατηγορίες χρηματοδοτικών εργαλείων ενεργειακής αναβάθμισης]{Ενδεικτικές κατηγορίες χρηματοδοτικών εργαλείων για ενεργειακές αναβαθμίσεις στην Ελλάδα.}
    \label{tab:financing-categories}
    \begin{tabularx}{\textwidth}{L{3.2cm}Y Y}
        \toprule
        \textbf{Κατηγορία} & \textbf{Περιγραφή} & \textbf{Ενδεικτικές πηγές πληροφορίας} \\
        \midrule
        Κρατικά προγράμματα επιδότησης & Επιχορηγήσεις ή κίνητρα για παρεμβάσεις εξοικονόμησης ενέργειας σε κατοικίες. & ΥΠΕΝ, επίσημες πύλες προγραμμάτων, Ταμείο Ανάκαμψης. \\
        Επιδοτούμενα δάνεια & Δανειακά προϊόντα με επιδότηση επιτοκίου ή ειδικούς όρους για ενεργειακές παρεμβάσεις. & Greece 2.0, τράπεζες, δημόσιες ανακοινώσεις. \\
        Τραπεζικά πράσινα προϊόντα & Επισκευαστικά ή καταναλωτικά δάνεια για πράσινες παρεμβάσεις, αγορά εξοπλισμού ή αναβάθμιση κατοικίας. & Ιστοσελίδες τραπεζών, ενημερωτικά φυλλάδια, όροι προϊόντων. \\
        Ιδιωτικά σχήματα και ESCO & Χρηματοδότηση ή υλοποίηση έργων μέσω συμβάσεων ενεργειακής απόδοσης. & Μητρώα επιχειρήσεων, συμβάσεις, οδηγοί ενεργειακών υπηρεσιών. \\
        \bottomrule
    \end{tabularx}
\end{table}

Η ποικιλία αυτών των εργαλείων καλύπτει διαφορετικές ανάγκες, αλλά αυξάνει και
την πολυπλοκότητα της πληροφόρησης. Κάθε εργαλείο έχει διαφορετικούς
δικαιούχους, προϋποθέσεις, επιλέξιμες παρεμβάσεις, όρια χρηματοδότησης και
διαδικασία αίτησης. Για αυτό η παρούσα εργασία περιορίζεται σκόπιμα σε
εργαλεία με σαφή ενεργειακή διάσταση σε κατοικίες και δεν επιχειρεί να
καταλογογραφήσει γενικά στεγαστικά ή επισκευαστικά προϊόντα χωρίς ενεργειακό
κριτήριο.

\subsection{Χρηματοδότηση ενεργειακών ανακαινίσεων και εμπόδια υλοποίησης}

Το αρχικό κόστος παραμένει ένα από τα βασικά εμπόδια για την υλοποίηση
ενεργειακών ανακαινίσεων. Ακόμη και όταν μια παρέμβαση μπορεί να αποδώσει σε
βάθος χρόνου, απαιτεί κεφάλαιο, τεχνική προετοιμασία και, σε πολυκατοικίες,
συντονισμό μεταξύ ιδιοκτητών με διαφορετικές δυνατότητες και προτεραιότητες.

Τα χρηματοδοτικά εργαλεία μειώνουν μέρος αυτού του εμποδίου μέσω επιχορηγήσεων,
δανείων, επιδότησης επιτοκίου, εγγυήσεων ή συμβάσεων ενεργειακής απόδοσης. Η
πρακτική τους αξία όμως εξαρτάται και από την ποιότητα της πληροφόρησης. Ο
πολίτης, ο μηχανικός ή ο σύμβουλος έργου χρειάζεται να γνωρίζει ποια εργαλεία
είναι ενεργά, ποιον αφορούν, ποιες παρεμβάσεις καλύπτουν, τι χρηματοδότηση
προβλέπουν και ποιες προθεσμίες ισχύουν. Επειδή οι όροι αυτοί αλλάζουν,
η πληροφορία δεν μπορεί να αντιμετωπίζεται ως στατικός κατάλογος, αλλά ως
δεδομένο που χρειάζεται τακτική επιστροφή στις πρωτογενείς πηγές.

\subsection{Κατακερματισμός πληροφορίας και ανάγκη υποστήριξης αποφάσεων}

Στην πράξη, οι πληροφορίες για χρηματοδοτικά εργαλεία ενεργειακής αναβάθμισης
είναι κατανεμημένες σε πολλαπλές πηγές. Επίσημες ανακοινώσεις δημοσιεύονται σε
ιστοσελίδες υπουργείων, οδηγοί εφαρμογής διατίθενται συχνά σε αρχεία PDF,
προγράμματα έχουν ξεχωριστές διαδικτυακές πύλες, ενώ οι τράπεζες παρουσιάζουν
τα δικά τους προϊόντα σε εμπορικές σελίδες με διαφορετική δομή και ορολογία. Η
δομή αυτού του προβλήματος παρουσιάζεται στο
Σχήμα~\ref{fig:information-fragmentation}.

\begin{figure}[!htbp]
    \centering
    \begin{adjustbox}{max width=\textwidth,center}
        \begin{tikzpicture}[
    source/.style={
        rectangle,
        rounded corners=3pt,
        draw=gray!60,
        fill=gray!8,
        align=center,
        text width=3.55cm,
        minimum height=1.05cm,
        font=\sffamily\small
    },
    hub/.style={
        circle,
        draw=orange!70!black,
        fill=orange!12,
        align=center,
        text width=3.2cm,
        minimum size=2.85cm,
        font=\sffamily\small\bfseries
    },
    problem/.style={
        rectangle,
        rounded corners=3pt,
        draw=red!55,
        fill=red!6,
        align=center,
        text width=12.5cm,
        minimum height=0.95cm,
        font=\sffamily\small
    },
    arrow/.style={-{Latex[length=2mm]}, thick, draw=gray!65}
]
\node[hub] (hub) at (0,0) {Χρηματοδοτική\\πληροφορία};

\node[source] (public) at (-4.9,2.25) {Δημόσιες πηγές\\ΥΠΕΝ, οδηγοί, προσκλήσεις};
\node[source] (portals) at (4.9,2.25) {Πύλες προγραμμάτων\\π.χ. Εξοικονομώ};
\node[source] (banks) at (-4.9,-2.25) {Τράπεζες\\πράσινα δάνεια};
\node[source] (docs) at (4.9,-2.25) {PDF και ενημερώσεις\\παραρτήματα, παρατάσεις};

\draw[arrow] (public) -- (hub);
\draw[arrow] (portals) -- (hub);
\draw[arrow] (banks) -- (hub);
\draw[arrow] (docs) -- (hub);

\node[problem] at (0,-4.05) {
Πρόβλημα: διαφορετική δομή, διαφορετική ορολογία, συχνές αλλαγές και δυσκολία σύγκρισης
};
\end{tikzpicture}

    \end{adjustbox}
    \caption[Κατακερματισμός πληροφορίας χρηματοδοτικών εργαλείων]{Κατακερματισμός πληροφορίας για χρηματοδοτικά εργαλεία ενεργειακής αναβάθμισης. Ιδία επεξεργασία.}
    \label{fig:information-fragmentation}
\end{figure}

Ο κατακερματισμός αυτός αυξάνει τον χρόνο αναζήτησης, δυσκολεύει τη σύγκριση
μεταξύ προγραμμάτων και κάνει απαιτητική τη συντήρηση ενημερωμένων δεδομένων.
Η χειροκίνητη καταγραφή μπορεί να λειτουργήσει για μικρό αριθμό πηγών, αλλά
γίνεται εύθραυστη όταν οι οδηγοί, οι προθεσμίες και οι τραπεζικοί όροι
μεταβάλλονται. Το κίνητρο της παρούσας διπλωματικής είναι ακριβώς αυτό το
ενδιάμεσο πρόβλημα: πώς μπορεί να αυτοματοποιηθεί μέρος της συλλογής και
τυποποίησης, χωρίς να χαθεί η σύνδεση με την επίσημη πηγή και χωρίς να
καταργηθεί ο ανθρώπινος έλεγχος.

\section{Σκοπός της διπλωματικής και προτεινόμενη λύση}

\subsection{Από την ετερογενή πληροφορία στη δομημένη γνώση}

Η πληροφορία που χρειάζεται ένας ενδιαφερόμενος για ένα χρηματοδοτικό εργαλείο
βρίσκεται συνήθως σε ελεύθερο κείμενο: ιστοσελίδες, οδηγούς εφαρμογής,
ανακοινώσεις, πίνακες ή αρχεία PDF. Για να αξιοποιηθεί από πληροφοριακό
σύστημα, πρέπει να μετατραπεί σε διακριτά πεδία, όπως όνομα προγράμματος,
φορέας, δικαιούχοι, επιλέξιμες παρεμβάσεις, ποσά χρηματοδότησης και
προθεσμίες.

Η τυποποίηση σε κοινό σχήμα επιτρέπει τη σύγκριση προγραμμάτων ακόμη και όταν
οι αρχικές πηγές έχουν διαφορετική δομή ή ορολογία. Δεν αντικαθιστά την
επίσημη πηγή, αλλά κάνει ορατό τι έχει εντοπιστεί, τι λείπει και ποια πεδία
χρειάζονται επιβεβαίωση. Η δυσκολία είναι ότι οι ίδιες έννοιες εμφανίζονται με
διαφορετικές διατυπώσεις: το «ανώτατο ποσό χρηματοδότησης» μπορεί να δηλώνεται
ως ύψος δανείου, επιλέξιμη δαπάνη, μέγιστος προϋπολογισμός ή ποσό
επιχορήγησης. Συνεπώς, η εξαγωγή δεν είναι απλή αναζήτηση λέξεων-κλειδιών, αλλά
εργασία κατανόησης συμφραζομένων.

\subsection{Μεγάλα γλωσσικά μοντέλα ως εργαλείο εξαγωγής και τυποποίησης πληροφορίας}

Τα μεγάλα γλωσσικά μοντέλα (\eng{Large Language Models} -- LLMs) είναι χρήσιμα
στην παρούσα εργασία επειδή μπορούν να επεξεργάζονται φυσική γλώσσα με βάση τα
συμφραζόμενα. Η αρχιτεκτονική τους και οι βασικές τεχνικές αξιοποίησης
παρουσιάζονται στο Κεφάλαιο~\ref{chap:theoretical-background}. Στην Εισαγωγή
αρκεί να αποσαφηνιστεί ο λειτουργικός τους ρόλος: χρησιμοποιούνται για να
διαβάζουν συγκεκριμένο περιεχόμενο, να αποφασίζουν αν μια πηγή είναι σχετική
και να εξάγουν πεδία σε προκαθορισμένη δομή.

Η χρήση τους πρέπει να παραμείνει περιορισμένη από την πηγή. Τα LLMs μπορούν να
παράγουν ανακριβείς ή μη τεκμηριωμένες απαντήσεις όταν καλούνται να
συμπληρώσουν κενά χωρίς επαρκές πλαίσιο. Για τον λόγο αυτό, η προτεινόμενη ροή
δεν ζητά από το μοντέλο να «γνωρίζει» τους όρους ενός προγράμματος, αλλά να
εργάζεται πάνω σε κείμενο που έχει συλλεχθεί από συγκεκριμένο URL. Με άλλα
λόγια, η προσέγγιση είναι τεκμηριωμένη στην πηγή (\eng{source-backed} ή
\eng{source-grounded}): κάθε εξαγωγή πρέπει να μπορεί να ελεγχθεί στην αρχική
πηγή της. Η ίδια λογική συνδέεται με προσεγγίσεις τεκμηριωμένης ανάκτησης γνώσης, όπως το
\eng{Retrieval-Augmented Generation}, όπου η παραγωγή απάντησης περιορίζεται
από εξωτερικό και ελέγξιμο περιεχόμενο \cite{lewis2020rag}.

Συνεπώς, το LLM αντιμετωπίζεται ως ενδιάμεσο εργαλείο μετασχηματισμού και όχι
ως τελικός κριτής. Η τελική πληροφορία πρέπει να συνοδεύεται από πηγή, χρονική
σήμανση και δυνατότητα ελέγχου πριν εμφανιστεί στον χρήστη της εφαρμογής.

\subsection{Ερευνητικό κενό, πεδίο εφαρμογής και συμβολή της εργασίας}

Παρότι υπάρχει εκτενής βιβλιογραφία για την ενεργειακή αποδοτικότητα, τα
χρηματοδοτικά εργαλεία και τα μεγάλα γλωσσικά μοντέλα, το ειδικό πρόβλημα της
αυτόματης εξαγωγής και τυποποίησης πληροφορίας για χρηματοδοτικά εργαλεία
ενεργειακής αναβάθμισης στην Ελλάδα παραμένει περιορισμένα διερευνημένο. Το
κενό δεν βρίσκεται μόνο στην τεχνική εξαγωγή κειμένου ούτε μόνο στην ανάλυση
πολιτικών, αλλά στη σύνδεσή τους σε μια ροή που μπορεί να ενημερώνεται, να
ελέγχεται και να χρησιμοποιείται από εφαρμογή.

Το πεδίο εφαρμογής περιορίζεται σε εργαλεία που σχετίζονται με ενεργειακές
αναβαθμίσεις κατοικιών ή οικιακές εφαρμογές ανανεώσιμων πηγών ενέργειας στην
Ελλάδα. Περιλαμβάνονται δημόσια προγράμματα, επιδοτούμενα ή πράσινα τραπεζικά
προϊόντα και συναφή εργαλεία με σαφή ενεργειακό κριτήριο. Αντίθετα, γενικά
στεγαστικά, επισκευαστικά ή καταναλωτικά προϊόντα χωρίς ενεργειακή διάσταση
δεν αποτελούν αντικείμενο της εργασίας. Η αντιστοίχιση του προβλήματος με την
τεχνική ανάγκη και τη συμβολή της εργασίας παρουσιάζεται στον
Πίνακα~\ref{tab:problem-contribution}.

\begin{table}[!htbp]
    \centering
    \caption[Πρόβλημα, τεχνική ανάγκη και συμβολή της εργασίας]{Συσχέτιση προβλήματος, τεχνικής ανάγκης και συμβολής της εργασίας.}
    \label{tab:problem-contribution}
    \begingroup
    \small
    \renewcommand{\arraystretch}{1.15}
    \begin{tabularx}{\textwidth}{L{3.2cm}Y Y}
        \toprule
        \textbf{Πρόβλημα} & \textbf{Τεχνική ανάγκη} & \textbf{Συμβολή εργασίας} \\
        \midrule
        Κατακερματισμένες πηγές & Στοχευμένη συλλογή περιεχομένου από επίσημες και τραπεζικές ιστοσελίδες. & Επαναλήψιμη ροή συλλογής και αποθήκευσης πηγών. \\
        Ανομοιογενής ορολογία & Αντιστοίχιση διαφορετικών διατυπώσεων σε κοινά πεδία. & Χρήση LLM για συνάφεια και εξαγωγή δομημένων δεδομένων. \\
        Συχνές αλλαγές προγραμμάτων & Ενημέρωση χωρίς χειροκίνητη επανακαταχώριση από την αρχή. & Χρονική σήμανση και σύνδεση κάθε εγγραφής με πρωτογενές URL. \\
        Κίνδυνος ανακριβών εξαγωγών & Έλεγχοι ποιότητας πριν από την αξιοποίηση της πληροφορίας. & Επικύρωση, κατάσταση ελέγχου και ανθρώπινη έγκριση. \\
        Δυσκολία αξιοποίησης από χρήστες & Καθαρή παρουσίαση μόνο πάνω σε εγκεκριμένα δεδομένα. & Διεπαφή αναζήτησης προγράμματος και υποβολής ερωτήσεων. \\
        \bottomrule
    \end{tabularx}
    \endgroup
\end{table}

Αντικείμενο της διπλωματικής είναι η ανάπτυξη μεθοδολογίας και εφαρμογής για
την αυτόματη εξαγωγή, τυποποίηση και αξιοποίηση δεδομένων που αφορούν τέτοια
εργαλεία. Η εργασία καλύπτει το στάδιο που προηγείται της τεχνικοοικονομικής
απόφασης: οργανώνει την πληροφορία ώστε ο χρήστης, ο μηχανικός ή ο
διαχειριστής να εντοπίζει γρήγορα τους βασικούς όρους ενός προγράμματος και να
επιστρέφει στην επίσημη πηγή για επιβεβαίωση.

Η προσέγγιση στηρίζεται σε τρεις παραδοχές: οι πρωτογενείς διαδικτυακές πηγές
αποτελούν το σημείο αναφοράς, η αυτόματη εξαγωγή λειτουργεί υποβοηθητικά και
η απουσία πληροφορίας πρέπει να δηλώνεται ως κενό πεδίο, όχι να καλύπτεται με
εικασία. Σε τεχνικό επίπεδο, η ροή ξεκινά από αξιόπιστες ή ρητά επιλεγμένες
πηγές, καθαρίζει το περιεχόμενο, αποφασίζει αν η πηγή είναι σχετική, εξάγει
πεδία σε κοινό JSON σχήμα και αποθηκεύει τα αποτελέσματα μαζί με στοιχεία
προέλευσης. Η πειραματική εκτέλεση υποστηρίζεται από τοπικά γλωσσικά μοντέλα
μέσω Ollama, ώστε η διαδικασία να μπορεί να επαναληφθεί με ελεγχόμενες
παραμέτρους \cite{ollama_docs_2026}.

Με βάση τα παραπάνω, η εργασία οργανώνεται γύρω από τρία ερευνητικά ερωτήματα:
(α) με ποια κριτήρια επιλέγονται και φιλτράρονται πηγές που αφορούν πράγματι
χρηματοδότηση ενεργειακών παρεμβάσεων, (β) πώς μετατρέπεται η ετερογενής
πληροφορία τους σε κοινό και ελέγξιμο σχήμα δεδομένων, και (γ) πώς
περιορίζεται ο κίνδυνος μη τεκμηριωμένης εξαγωγής από το LLM μέσω
ιχνηλασιμότητας, ελέγχων και ανθρώπινης εποπτείας.

Οι βασικοί στόχοι της εργασίας είναι:

\begin{enumerate}
    \item Ο σαφής ορισμός του πεδίου ενδιαφέροντος και των κατηγοριών πηγών.
    \item Ο σχεδιασμός ροής για συλλογή περιεχομένου, απόφαση συνάφειας,
    εξαγωγή πεδίων και αποθήκευση σε δομημένο σχήμα.
    \item Η αξιοποίηση LLMs για αναγνώριση σχετικών πηγών και εξαγωγή πεδίων
    όπως δικαιούχοι, παρεμβάσεις, ποσά, επιτόκια, προθεσμίες και διαδικασία
    αίτησης.
    \item Η ενσωμάτωση ελέγχου ποιότητας, ανθρώπινης εποπτείας και αξιολόγησης
    ως προς συνάφεια, πληρότητα, συνέπεια και τεκμηρίωση.
\end{enumerate}

Η συνολική ροή της προτεινόμενης προσέγγισης, από την είσοδο των πηγών μέχρι
τη χρήση των εγκεκριμένων δεδομένων στην εφαρμογή, παρουσιάζεται στο
Σχήμα~\ref{fig:pipeline-concept}.

\begin{figure}[!htbp]
    \centering
    \begin{adjustbox}{max width=\textwidth,center}
        \begin{tikzpicture}[
    card/.style={
        rectangle,
        rounded corners=3pt,
        draw=blue!55,
        fill=blue!6,
        align=left,
        text width=4.95cm,
        minimum height=0.92cm,
        inner xsep=7pt,
        font=\sffamily\small
    },
    outcome/.style={
        rectangle,
        rounded corners=3pt,
        draw=teal!55!black,
        fill=teal!7,
        align=left,
        text width=4.95cm,
        minimum height=0.92cm,
        inner xsep=7pt,
        font=\sffamily\small
    },
    arrow/.style={-{Latex[length=2mm]}, thick, draw=gray!65}
]
\node[card] (s1) at (0,0) {\textbf{1.} Διαδικτυακές πηγές};
\node[card, below=0.36cm of s1] (s2) {\textbf{2.} Συλλογή και καθαρισμός περιεχομένου};
\node[card, below=0.36cm of s2] (s3) {\textbf{3.} Απόφαση συνάφειας με LLM};
\node[card, below=0.36cm of s3] (s4) {\textbf{4.} Εξαγωγή σε δομημένο JSON};

\node[outcome, right=1.05cm of s1] (s5) {\textbf{5.} Έλεγχος ποιότητας και τεκμηρίωσης};
\node[outcome, below=0.36cm of s5] (s6) {\textbf{6.} Δομημένη αποθήκευση JSON};
\node[outcome, below=0.36cm of s6] (s7) {\textbf{7.} Web εφαρμογή διαχειριστή};
\node[outcome, below=0.36cm of s7] (s8) {\textbf{8.} Προβολή και ερωτήσεις χρήστη};

\draw[arrow] (s1) -- (s2);
\draw[arrow] (s2) -- (s3);
\draw[arrow] (s3) -- (s4);
\draw[arrow] (s4.east) -- ++(0.52,0) |- (s5.west);
\draw[arrow] (s5) -- (s6);
\draw[arrow] (s6) -- (s7);
\draw[arrow] (s7) -- (s8);
\end{tikzpicture}

    \end{adjustbox}
    \caption[Εννοιολογική ροή της προτεινόμενης προσέγγισης]{Εννοιολογική ροή της προτεινόμενης προσέγγισης: από τις ετερογενείς πηγές πληροφορίας στη δομημένη γνώση και στη διαδικτυακή εφαρμογή. Ιδία επεξεργασία.}
    \label{fig:pipeline-concept}
\end{figure}

Η συμβολή της εργασίας βρίσκεται στον συνδυασμό ενεργειακής πολιτικής,
χρηματοδοτικής πληροφορίας και τεχνικών επεξεργασίας φυσικής γλώσσας σε μια
ενιαία, ελέγξιμη ροή. Η συγκέντρωση της πληροφορίας μειώνει την αναζήτηση σε
πολλές πηγές, ενώ η σύνδεση κάθε εγγραφής με την πηγή της υποστηρίζει έλεγχο,
σύγκριση και μελλοντική επικαιροποίηση.

\section{Οργάνωση Τόμου}

Η συνέχεια της διπλωματικής εργασίας οργανώνεται ως εξής. Στο
Κεφάλαιο~\ref{chap:theoretical-background} παρουσιάζεται το θεωρητικό υπόβαθρο
της εργασίας: ενεργειακή αναβάθμιση κτιρίων, χρηματοδοτικά εργαλεία, μεγάλα
γλωσσικά μοντέλα και βασικοί περιορισμοί αξιοπιστίας, κόστους και ανθρώπινης
εποπτείας.

Στο Κεφάλαιο~\ref{chap:methodology} περιγράφεται η μεθοδολογία. Παρουσιάζονται
το πεδίο εφαρμογής, οι απαιτήσεις, η αρχιτεκτονική και τα βασικά βήματα με τα
οποία συνδυάζονται η διαδικτυακή συλλογή δεδομένων και τα LLMs, χωρίς να
αναλύονται ακόμη οι λεπτομέρειες κώδικα.

Στο κεφάλαιο της υλοποίησης αναλύονται τα επιμέρους τμήματα της εφαρμογής,
όπως η επιλογή και επεξεργασία συνδέσμων, η εξαγωγή περιεχομένου, η χρήση
μοντέλων για απόφαση συνάφειας, η παραγωγή δομημένων δεδομένων, οι μηχανισμοί
ελέγχου και η αποθήκευση των αποτελεσμάτων. Επίσης, τεκμηριώνονται οι τεχνικές
επιλογές, τα μοντέλα που εξετάστηκαν και οι βασικές παράμετροι χρήσης τους.

Στο κεφάλαιο των αποτελεσμάτων παρουσιάζεται η αξιολόγηση της προσέγγισης,
κυρίως ως προς την απόδοση των LLMs και την ποιότητα των παραγόμενων δομημένων
δεδομένων. Περιλαμβάνονται έλεγχοι συνέπειας ανάμεσα στις απαντήσεις και στα
εξαγόμενα πεδία, επαναληψιμότητα απαντήσεων, σημασιολογική σύγκριση με
αναμενόμενες απαντήσεις, έλεγχος τεκμηρίωσης στο αρχικό HTML και γραφικές
απεικονίσεις των βασικών μετρικών.

Στο κεφάλαιο της διαδικτυακής εφαρμογής αναλύονται η ροή δεδομένων, η
διάκριση ρόλων διαχειριστή και τελικού χρήστη, τα βήματα χρήσης και ενδεικτικά
παραδείγματα αλληλεπίδρασης με τα δεδομένα και το σύστημα ερωταπαντήσεων.
Ιδιαίτερη έμφαση δίνεται στη ροή διαχειριστή, όπου μπορούν να δοθούν URLs προς
επεξεργασία, να παρακολουθηθεί η εκτέλεση εργασίας και να ελεγχθούν τα
παραγόμενα JSON αρχεία πριν αξιοποιηθούν από τον τελικό χρήστη.

Τέλος, στο κεφάλαιο των συμπερασμάτων συνοψίζονται τα βασικά ευρήματα της
εργασίας, σχολιάζεται η καινοτομία και η πρακτική χρησιμότητα της προσέγγισης
και προτείνονται μελλοντικές κατευθύνσεις βελτίωσης, όπως η επέκταση σε
περισσότερες πηγές, η ενσωμάτωση μηχανισμών αυτόματης παρακολούθησης αλλαγών
και η βελτίωση των διαδικασιών αξιολόγησης και επαλήθευσης.
